% Part: many-valued-logic
% Chapter: three-valued-logics
% Section: goedel

\documentclass[../../../include/open-logic-section]{subfiles}

\begin{document}

\olfileid{mvl}{thr}{god}

\olsection{منطق‌های گودل}

کورت گودل دنباله‌ای از منطق‌های~$n$ارزشی معرفی کرد که هریک همهٔ
!!{formula}های معتبر در منطق شهودگرایانه را در بر دارند و خود در منطق
کلاسیک گنجانده می‌شوند. نخستین مورد جالب چنین است:

\begin{defn}\ollabel{defn:goedel}
\emph{منطق سه‌ارزشی گودل}~$\LogGod$ با ماتریس زیر تعریف می‌شود:
\begin{enumerate}
  \item زبان گزاره‌ای استاندارد~$\Lang L_0$ با~$\lfalse$، $\lnot$،
  $\land$، $\lor$ و~$\lif$.
  \item مجموعهٔ ارزش‌های صدق~$V = \{\True, \Undef, \False\}$.
  \item $\True$ تنها ارزش ممتاز است؛ یعنی~$V^+ = \{\True\}$.
  \item برای~$\lfalse$، داریم~$\tf{\lfalse} = \False$. توابع صدق
  رابط‌های دیگر با جدول‌های زیر داده می‌شوند:
  \begin{center}
    \begin{tabular}{c|c} 
      $\tf{\lnot}[\LogGod]$ & \\ 
      \hline  
      $\True$ & $\False$ \\ 
      $\Undef$ & $\False$ \\
      $\False$ & $\True$ 
    \end{tabular}
    \quad
    \begin{tabular}{c|ccc} 
      $\tf{\land}[\LogGod]$ & $\True$ & $\Undef$ & $\False$ \\ 
      \hline 
      $\True$ & $\True$ & $\Undef$ & $\False$ \\ 
      $\Undef$ & $\Undef$ & $\Undef$ & $\False$\\ 
      $\False$ & $\False$ & $\False$ & $\False$ 
    \end{tabular}
    \\[2ex]
    \begin{tabular}{c|ccc} 
      $\tf{\lor}[\LogGod]$ & $\True$ & $\Undef$ & $\False$ \\ 
      \hline 
      $\True$ & $\True$ & $\True$ & $\True$ \\ 
      $\Undef$ & $\True$ & $\Undef$ & $\Undef$ \\
      $\False$ & $\True$ & $\Undef$ & $\False$ 
    \end{tabular}
    \quad
    \begin{tabular}{c|ccc} 
      $\tf{\lif}[\LogGod]$ & $\True$ & $\Undef$ & $\False$ \\ 
      \hline 
      $\True$ & $\True$ & $\Undef$ & $\False$ \\ 
      $\Undef$ & $\True$ & $\True$ & $\False$  \\ 
      $\False$ & $\True$ & $\True$ & $\True$ 
    \end{tabular}
  \end{center} 
\end{enumerate}
\end{defn}

جدول‌های صدق~$\land$ و~$\lor$ همان جدول‌های منطق \L ukasiewicz و
منطق قوی کلینی‌اند، اما جدول‌های صدق~$\lnot$ و~$\lif$ در هریک متفاوت
است. در منطق گودل، $\tf{\lnot}(\Undef) = \False$. برخلاف منطق
\L ukasiewicz و منطق کلینی، $\tf{\lif}(\Undef, \False) = \False$؛ و
برخلاف منطق کلینی، اما مانند منطق \L ukasiewicz،
$\tf{\lif}(\Undef, \Undef) = \True$.

چنان‌که پیوند اشاره‌شده با منطق شهودگرایانه نشان می‌دهد، $\LogGod[3]$
به منطق شهودگرایانه نزدیک است. همهٔ حقایق شهودگرایانه در~$\LogGod[3]$
همان‌گویی‌اند و بسیاری از همان‌گویی‌های کلاسیک که از دید شهودگرایانه
معتبر نیستند، در~$\LogGod[3]$ نیز همان‌گویی نیستند. برای نمونه، عبارات
زیر همان‌گویی نیستند:
\begin{align*}
  & p \lor \lnot p && (p \lif q) \lif (\lnot p \lor q) \\
  & \lnot\lnot p \lif p && \lnot(\lnot p \land \lnot q) \lif (p \lor q) \\
  & ((p \lif q) \lif p) \lif p && \lnot(p \lif q) \lif (p \land \lnot q)
\end{align*}
بااین‌حال، هر همان‌گویی~$\LogGod[3]$ نیز لزوماً از دید شهودگرایانه
معتبر نیست؛ برای نمونه~$\lnot\lnot p \lor \lnot p$ یا
$(p \lif q) \lor (q \lif p)$.

\begin{prob}
  جدول‌های صدقی بدهید که نشان دهند عبارات زیر همان‌گویی‌های
  $\LogGod[3]$ هستند:
  \begin{align*}
    & \lnot\lnot p \lor \lnot p\\
    & (p \lif q) \lor (q \lif p) \\
    & \lnot(p \land q) \lif (\lnot p \lor \lnot q) \\
    & (p \lif q) \lor (q \lif r) \lor (r \lif s)
  \end{align*}
\end{prob}

\begin{prob}
  جدول‌های صدقی بدهید که نشان دهند عبارات زیر همان‌گویی‌های
  $\LogGod[3]$ نیستند:
  \begin{align*}
    & (p \lif q) \lif (\lnot p \lor q) \\
    & \lnot(\lnot p \land \lnot q) \lif (p \lor q) \\
    & ((p \lif q) \lif p) \lif p \\
    & \lnot(p \lif q) \lif (p \land \lnot q)
  \end{align*}
\end{prob}

\begin{prob}
  کدام‌یک از روابط زیر در منطق گودل برقرارند؟ برای هریک جدول صدق بدهید.
  \begin{enumerate}
    \item $p, p \lif q \Entails q$
    \item $p \lor q, \lnot p \Entails q$
    \item $p \land q \Entails p$
    \item $p \Entails p \land p$
    \item $p \Entails p \lor q$
  \end{enumerate}
\end{prob}

\end{document}
